\section{Setting and definitions}\label{sec:setting}

\paragraph{Setting.} $G=(V,E)$ is a graph in the sense of the textbook:
a set $V=\{n_1,\dots,n_N\}$ of vertices, $N\ge 1$, and a set $E$ of edges
\cite[Definition~1.1.1, p.~9]{petrosyan2015graphs}. It is finite, simple
(no edge from a vertex to itself, at most one edge between two vertices) and
undirected (the edge joining $u$ and $v$, written $uv$, is the same as $vu$)
\cite{userstatements}. The vertices $u$ and $v$ are \emph{adjacent} if
$uv\in E$ \cite[Definition~1.1.4, p.~11]{petrosyan2015graphs}. The indices
$i,j,r$ run over $1,\dots,N$; $k\ge 0$ is an integer.

\plainwords{} A pond has islands, and some pairs of islands have a bridge between them. Draw a square table with one row and one column for each island. Put a one in a cell when the two islands are next door, and a zero in every other cell.

\begin{definition}[Adjacency matrix]\label{def:matrix}
The \emph{adjacency matrix} of $G$ is the $N\times N$ matrix $A$ with
$A_{ij}=1$ if $n_i$ and $n_j$ are adjacent and $A_{ij}=0$ otherwise
\cite[p.~11]{petrosyan2015graphs}; the board states the same rule in other
letters \cite{whiteboard1,whiteboard2}.
\end{definition}
Its entries are $0$ or $1$. Matrix powers are the usual ones: $(XY)_{ij}=\sum_{r=1}^{N}X_{ir}Y_{rj}$,
$A^0=I$ (the identity matrix) and $A^{k+1}=A^kA$. Every entry of
$A^k$ is a non-negative integer, since sums and products of non-negative
integers are non-negative integers (\emph{used in} \cref{thm:sum}).

\plainwords{} A frog sits on an island and hops over a bridge to a next-door island. It may hop back, or land on the same island again. The islands it visits, in order, form a trip. A frog that never hops makes a trip with no hops.

\begin{definition}[Walk]\label{def:walk}
A \emph{walk} from $u_0$ to $u_k$ is a sequence of vertices
$u_0u_1\cdots u_k$ in which each vertex is adjacent to the one before it; its
\emph{length} is $k$, the number of edges used, each counted as often as it
is used. Vertices and edges may repeat. This is the textbook's Definition~1.2.10
\cite[p.~21]{petrosyan2015graphs}, which lists edges as well; in a simple graph
the vertices fix the edges. This paper also allows $k=0$: a single vertex is a walk of length $0$ from itself to
itself. This is the paper's own convention; the textbook's usage assumes it
\cite[pp.~23, 49]{petrosyan2015graphs}.
\end{definition}

\plainwords{} A fresh-island trip lands on no island twice. So it never crosses the same bridge twice either. Every fresh-island trip is a trip, but a trip may land on an island twice.

\begin{definition}[Path]\label{def:path}
A \emph{path} is a walk whose vertices are pairwise distinct. This is the
course convention \cite{userstatements}. It is the textbook's \emph{simple
path} \cite[Definition~1.2.13, p.~22]{petrosyan2015graphs}; the textbook's
own word ``path'' means a walk with no repeated edge, and this paper never
uses the word in that sense.
\end{definition}

\plainwords{} A pond is in one piece when the frog can get from any island to any other. The way may cross many bridges. It is enough to ask for a fresh-island trip between each two islands.

\begin{definition}[Connected]\label{def:connected}
$G$ is \emph{connected} if for any two vertices $u,v$ there is a path from
$u$ to $v$. The board states this
\cite{whiteboard1,whiteboard2}; the textbook's Definition~2.1.1
\cite[p.~48]{petrosyan2015graphs} is the same with the textbook's own ``path''. A vertex is joined to itself by the path
of length $0$, so only pairs of different vertices matter.
\end{definition}

\paragraph{Board graph.} The board's matrix has rows $0\,1\,0$, $1\,0\,1$,
$0\,1\,0$ \cite{whiteboard1,whiteboard2}:
\[
A=\begin{pmatrix}0&1&0\\1&0&1\\0&1&0\end{pmatrix}.
\]
It encodes the graph with edges $n_1n_2$ and $n_2n_3$, which is the example
graph throughout. The board's drawing shows a third line that the matrix
does not encode; the matrix defines the example.
